\chapter{复数}\label{ch:g12:complex}

方程 $x^2 = -1$ 没有实数解。将 $\R$ 扩充，添加一个新数 $\iu$，
其平方为 $-1$，得到复数域 $\C$，
其中\emph{每一个}多项式方程都有解 ---
且其运算刻画了平面的几何：平移、旋转
与伸缩分别对应为加法与乘法。

\section{代数形式}

\begin{definition}[复数]\label{def:g12:complex:def}
所有\emph{复数}\index{复数}组成的集合为
\[
\C = \{a + \iu b : a, b \in \R\},
\]
其中 $\iu$ 为满足唯一规则 $\iu^2 = -1$ 的符号；加法与乘法
按 $\iu$ 的实系数多项式的方法进行，并将 $\iu^2$ 化为 $-1$。
实数 $a$ 与 $b$ 分别称为 \emph{实部}\index{实部}
$\Rea(z)$ 与 \emph{虚部}\index{虚部} $\Ima(z)$，它们是 $z = a + \iu b$ 的；这一表达式
称为 $z$ 的\emph{代数形式}\index{代数形式}，且唯一：
$a + \iu b = a' + \iu b'$ 当且仅当 $a = a'$ 且 $b = b'$。
\end{definition}

\begin{example}
$(2 + 3\iu)(1 - \iu) = 2 - 2\iu + 3\iu - 3\iu^2 = 5 + \iu$.
\end{example}

\begin{definition}[共轭复数]\label{def:g12:complex:conjugate}
$z = a + \iu b$ 的\emph{共轭复数}\index{共轭复数}为
$\conj{z} = a - \iu b$.
\end{definition}

\begin{proposition}[共轭运算的性质]\label{prop:g12:complex:conjugate}
对所有 $z, w \in \C$：
\[
\conj{z + w} = \conj z + \conj w, \quad
\conj{zw} = \conj z\,\conj w, \quad
\conj{\conj z} = z, \quad
z + \conj z = 2\Rea(z), \quad
z - \conj z = 2\iu\,\Ima(z),
\]
且 $z\conj z = a^2 + b^2 \geq 0$，其中 $z = a + \iu b$。此外，$z \in \R$
当且仅当 $z = \conj z$。
\end{proposition}

\begin{proof}
以上均可由代数形式直接计算得到；例如
$z\conj z = (a + \iu b)(a - \iu b) = a^2 - (\iu b)^2 = a^2 + b^2$.
\end{proof}

\begin{method}[复数的除法]\label{met:g12:complex:division}
要将商写成代数形式，只需将分子与分母同乘以
分母的共轭复数：
\[
\frac{1}{2 + \iu} = \frac{2 - \iu}{(2+\iu)(2-\iu)} = \frac{2 - \iu}{5}
= \frac25 - \frac15\iu .
\]
因此每个非零复数都有倒数，从而 $\C$ 是一个
\emph{域}。
\end{method}

\section{复平面、模与辐角}

\begin{definition}[对应复数、模]\label{def:g12:complex:modulus}
在标准正交坐标系的平面中，点 $M(a, b)$ 是
$z = a + \iu b$ 的\emph{像}，且 $z$ 为 $M$ 的\emph{对应复数}\index{对应复数}。
$z$ 的\emph{模}\index{模}是点 $M$ 到原点的距离：
\[
\abs{z} = \sqrt{a^2 + b^2} = \sqrt{z\conj z}.
\]
\end{definition}

\begin{omfigure}
\begin{tikzpicture}[scale=1.5]
  \draw[->] (-0.7,0) -- (2.9,0) node[right] {$x$};
  \draw[->] (0,-1.9) -- (0,2.1) node[above] {$y$};
  \draw[dashed] (2,1.5) -- (2,0) node[below right] {$a$};
  \draw[dashed] (2,1.5) -- (0,1.5) node[left] {$b$};
  \draw[dashed] (2,-1.5) -- (2,0);
  \draw[-Stealth,omDef,thick] (0,0) -- (2,1.5)
    node[midway, above left, font=\small] {$\abs z$};
  \draw[-Stealth,omThm,thick] (0,0) -- (2,-1.5);
  \fill (2,1.5) circle (1pt) node[above right] {$M(z = a + \iu b)$};
  \fill (2,-1.5) circle (1pt) node[below right] {$\conj z = a - \iu b$};
\end{tikzpicture}

{\small 点 $M(a,b)$ 的对应复数为 $z$：模 $\abs z$ 是
线段 $OM$ 的长度，共轭复数 $\conj z$ 是 $z$ 关于
实轴的对称点。}
\end{omfigure}

\begin{proposition}[模的性质]\label{prop:g12:complex:modulus}
对 $z, w \in \C$：
\[
\abs{zw} = \abs z\,\abs w, \qquad
\abs{\tfrac zw} = \tfrac{\abs z}{\abs w}\ (w \neq 0), \qquad
\abs{\conj z} = \abs z, \qquad
\abs{z + w} \leq \abs{z} + \abs{w} \ \text{（三角不等式）}.
\]
此外，$\abs{z_B - z_A}$ 是对应复数为
$z_A$ 与 $z_B$ 的两点间的距离。
\end{proposition}

\begin{proof}
$\abs{zw}^2 = zw\,\conj{zw} = z\conj z\, w \conj w = \abs z^2 \abs w^2$
即得第一式；第二式与第三式类似。对于三角
不等式：
\[
\abs{z+w}^2 = (z+w)(\conj z + \conj w)
= \abs z^2 + \abs w^2 + 2\Rea(z \conj w)
\leq \abs z^2 + \abs w^2 + 2\abs{z\conj w}
= \bigl(\abs z + \abs w\bigr)^2,
\]
利用 $\Rea(u) \leq \abs u$。距离的结论即为将勾股
公式用于 $z_B - z_A$ 的坐标。
\end{proof}

\begin{definition}[辐角、三角形式与指数形式]\label{def:g12:complex:argument}
设 $z \neq 0$。$z$ 的\emph{辐角}\index{辐角}，记作
$\Arg(z)$，是任一度量 $\theta$，它表示从正实轴到射线 $OM$ 所成的角；
它在相差 $2k\pi$ 的意义下确定。记 $r = \abs z$，
\[
z = r(\cos\theta + \iu\sin\theta)
\qquad\text{（三角形式）}.
\]
引入记号\index{指数形式}
\[
\eu^{\iu\theta} = \cos\theta + \iu\sin\theta,
\]
于是得到\emph{指数形式} $z = r\,\eu^{\iu\theta}$。
\end{definition}

\begin{omfigure}
\begin{tikzpicture}[scale=1.5]
  \draw[->] (-0.7,0) -- (2.4,0) node[right] {$x$};
  \draw[->] (0,-0.5) -- (0,1.9) node[above] {$y$};
  \draw[gray, dashed] (1.9,0) arc (0:75:1.9);
  \draw[-Stealth,omDef,thick] (0,0) -- (40:1.9)
    node[pos=0.62, below right=1pt, font=\small] {$r = \abs z$};
  \fill (40:1.9) circle (1pt) node[above right] {$z = r\,\eu^{\iu\theta}$};
  \draw[->,omThm,thick] (0.45,0) arc (0:40:0.45);
  \node[omThm] at (20:0.65) {$\theta$};
\end{tikzpicture}

{\small 指数形式：$z$ 由它到原点的距离 $r$ 与
从正实轴量起的角 $\theta$ 唯一确定。}
\end{omfigure}

\begin{theorem}\label{thm:g12:complex:expform}
对所有 $\theta, \theta' \in \R$：
\[
\eu^{\iu\theta}\,\eu^{\iu\theta'} = \eu^{\iu(\theta + \theta')},
\qquad
\conj{\eu^{\iu\theta}} = \eu^{-\iu\theta},
\qquad
\abs{\eu^{\iu\theta}} = 1 .
\]
由此，对非零的 $z = r\eu^{\iu\theta}$ 与
$z' = r'\eu^{\iu\theta'}$，有：
\[
zz' = rr'\,\eu^{\iu(\theta+\theta')},
\qquad
\frac{z}{z'} = \frac{r}{r'}\,\eu^{\iu(\theta - \theta')} :
\]
\emph{模相乘，辐角相加。}
\end{theorem}

\begin{proof}
第一个恒等式\emph{即是}两个加法公式
（\cref{prop:g12:trigo:addition}）：
\[
(\cos\theta + \iu\sin\theta)(\cos\theta' + \iu\sin\theta')
= (\cos\theta\cos\theta' - \sin\theta\sin\theta')
+ \iu(\sin\theta\cos\theta' + \cos\theta\sin\theta').
\]
其余部分由直接计算可得。
\end{proof}

\begin{corollary}[棣莫弗公式]\label{cor:g12:complex:demoivre}
\index{棣莫弗公式}
对 $\theta \in \R$ 与 $n \in \Z$，有
$\bigl(\cos\theta + \iu\sin\theta\bigr)^n = \cos n\theta + \iu \sin n\theta$。
\end{corollary}

\begin{proof}
对 $n \geq 0$ 用数学归纳法，并利用\cref{thm:g12:complex:expform}；
再通过取逆，将结论推广到 $n < 0$ 的情形。
\end{proof}

\begin{method}[两种形式之间的互化]\label{met:g12:complex:forms}
由代数形式化为指数形式：计算 $r = \abs z$，再求角 $\theta$，使
$\cos\theta = \frac ar$，$\sin\theta = \frac br$（使用反余弦前，
须先定出正确的象限）。由指数形式化为代数形式：展开
$r\cos\theta + \iu\, r\sin\theta$。求和用代数形式，
乘积、幂与商用指数形式。
\end{method}

\begin{example}
取 $z = 1 + \iu$，则 $\abs z = \sqrt2$，且 $\cos\theta = \sin\theta =
\frac{1}{\sqrt2}$，故 $\theta = \frac\pi4$，且 $z = \sqrt2\,\eu^{\iu\pi/4}$。
因此 $z^{20} = 2^{10}\,\eu^{5\iu\pi} = -1024$。
\end{example}

\section{二次方程与单位根}

\begin{theorem}[实系数二次方程]\label{thm:g12:complex:quadratic}
设 $a, b, c \in \R$，$a \neq 0$，$\Delta = b^2 - 4ac$。方程
$az^2 + bz + c = 0$ 在 $\C$ 中的解为：
\[
\begin{aligned}
&\Delta > 0:\ z = \frac{-b \pm \sqrt{\Delta}}{2a} \ (\text{两个实根}); \qquad
\Delta = 0:\ z = \frac{-b}{2a} \ (\text{二重根}); \\
&\Delta < 0:\ z = \frac{-b \pm \iu\sqrt{-\Delta}}{2a} \ (\text{两个共轭复根}).
\end{aligned}
\]
\end{theorem}

\begin{proof}
配方：
$az^2 + bz + c = a\left(\left(z + \frac{b}{2a}\right)^2 -
\frac{\Delta}{4a^2}\right)$。若 $\Delta < 0$，则
$\frac{\Delta}{4a^2} = \left(\iu\frac{\sqrt{-\Delta}}{2a}\right)^2$，于是
由平方差公式照常分解。
\end{proof}

\begin{definition}[单位根]\label{def:g12:complex:rootsofunity}
对 $n \geq 1$，\emph{$n$ 次单位根}\index{单位根}是
方程 $z^n = 1$ 的解。
\end{definition}

\begin{theorem}\label{thm:g12:complex:rootsofunity}
方程 $z^n = 1$ 恰有 $n$ 个解：
\[
\omega_k = \eu^{2\iu k\pi/n}, \qquad k = 0, 1, \dots, n-1 .
\]
它们的像为单位圆内接正 $n$ 边形的顶点，
且当 $n \geq 2$ 时，它们的和为 $0$。
\end{theorem}

\begin{proof}
每个 $\omega_k$ 满足 $\omega_k^n = \eu^{2\iu k\pi} = 1$，且它们
两两互异，因为它们的辐角属于 $\intco{0}{2\pi}$。
反之，若 $z^n = 1$，则 $\abs z^n = 1$，于是 $\abs z = 1$（正
实数），且 $z = \eu^{\iu\theta}$，其中 $n\theta \equiv 0 \pmod{2\pi}$，
\emph{即} $\theta = \frac{2k\pi}{n}$；将 $k$ 模 $n$ 约化，
便得某个 $\omega_k$。各顶点等角间隔 $\frac{2\pi}{n}$，
故为正 $n$ 边形。最后，取 $\omega = \omega_1$，
且 $\omega \neq 1$，由等比求和公式得
\[
\sum_{k=0}^{n-1}\omega_k = \sum_{k=0}^{n-1}\omega^k
= \frac{\omega^n - 1}{\omega - 1} = 0 . \qedhere
\]
\end{proof}

\begin{omfigure}
\begin{tikzpicture}[scale=1.7]
  \draw[->] (-1.3,0) -- (1.45,0) node[right] {$x$};
  \draw[->] (0,-1.3) -- (0,1.4) node[above] {$y$};
  \draw[gray] (0,0) circle (1);
  \draw[omDef] (0:1) -- (72:1) -- (144:1) -- (216:1) -- (288:1) -- cycle;
  \draw[->,omThm,thick] (0.3,0) arc (0:72:0.3);
  \node[omThm, font=\small] at (36:0.52) {$\frac{2\pi}{5}$};
  \fill (0:1) circle (1pt) node[below right] {$1$};
  \fill (72:1) circle (1pt) node[above right] {$\omega$};
  \fill (144:1) circle (1pt) node[above left] {$\omega^2$};
  \fill (216:1) circle (1pt) node[below left] {$\omega^3$};
  \fill (288:1) circle (1pt) node[below right] {$\omega^4$};
\end{tikzpicture}

{\small 五次单位根，$\omega = \eu^{2\iu\pi/5}$：
为单位圆内接正五边形的顶点，
（见\cref{exo:g12:complex:10}）。}
\end{omfigure}

\begin{example}
立方单位根为 $1$、$j = \eu^{2\iu\pi/3} =
-\frac12 + \iu\frac{\sqrt3}{2}$ 及 $\conj j = j^2$，满足
$1 + j + j^2 = 0$。
\end{example}

\section{复数与平面几何}

\begin{proposition}[几何对应]\label{prop:g12:complex:geometry}
设点 $A, B, C, D$ 的对应复数分别为 $a, b, c, d$（$a \neq b$，
$c \neq d$）。
\begin{enumerate}
  \item 由对应复数为 $t$ 的向量确定的平移为 $z \mapsto z + t$。
  \item 中心为 $\omega$（对应复数）、角为 $\theta$ 的旋转为
        $z \mapsto \omega + \eu^{\iu\theta}(z - \omega)$。
  \item 中心为 $\omega$、比为
        $k \in \R^*$ 的伸缩（位似）变换为 $z \mapsto \omega + k(z - \omega)$。
  \item $\dfrac{d - c}{b - a}$ 的模为 $\dfrac{CD}{AB}$，辐角为
        向量 $\vect{AB}$ 与 $\vect{CD}$ 之间的夹角。特别地，
        $AB \perp CD$ 当且仅当 $\frac{d-c}{b-a}$ 为纯虚数，
        且 $A, B, C$ 共线（$C \ne A$，$C\ne B$），当且仅当
        $\frac{c-a}{b-a} \in \R$。
\end{enumerate}
\end{proposition}

\begin{proof}
(1)、(3) 只是重述坐标公式。对 (2)，映射
$w \mapsto \eu^{\iu\theta} w$ 将模乘以 $1$，并将辐角加上 $\theta$，
即绕原点旋转角 $\theta$；再经把 $\omega$ 映为 $0$ 的平移作共轭，
即得一般中心的情形。对
(4)，商的模与辐角分别作差
（见\cref{thm:g12:complex:expform}）。
\end{proof}

\section{练习}

\begin{exercise}[$\star$]\label{exo:g12:complex:1}
将下列复数化为代数形式：
$(3 - 2\iu)(1 + 4\iu)$, \quad $\dfrac{2 + \iu}{1 - \iu}$, \quad
$\iu^{2027}$.
\end{exercise}

\begin{exercise}[$\star$]\label{exo:g12:complex:2}
在 $\C$ 中解方程：(a) $z^2 - 4z + 13 = 0$；(b) $z^2 + z + 1 = 0$。
在两种情形中验证两解互为共轭复数，并求它们的乘积。
\end{exercise}

\begin{exercise}[$\star$]\label{exo:g12:complex:3}
将 $z_1 = -1 + \iu$、$z_2 = \sqrt3 - \iu$ 化为指数形式，
并先以指数形式、再以代数形式计算 $\dfrac{z_1}{z_2}$，由此推出
$\cos\dfrac{11\pi}{12}$ 的精确值。
\end{exercise}

\begin{exercise}[$\star\star$]\label{exo:g12:complex:4}
描述满足下列条件的点 $M$（对应复数为 $z$）所构成的集合：
(a) $\abs{z - 2} = \abs{z + \iu}$；\quad
(b) $\abs{z - 1 - \iu} = 3$；\quad
(c) $\dfrac{z - 1}{z + 1}$ 为纯虚数。
\end{exercise}

\begin{exercise}[$\star\star$]\label{exo:g12:complex:5}
利用棣莫弗公式与二项式定理，将 $\cos 3\theta$
表示为 $\cos\theta$ 的多项式，并将 $\sin 3\theta$ 用
$\sin\theta$ 表示。
\end{exercise}

\begin{exercise}[$\star\star$]\label{exo:g12:complex:6}
\emph{（线性化。）}利用
$\cos\theta = \frac{\eu^{\iu\theta} + \eu^{-\iu\theta}}{2}$，将
$\cos^3\theta$ 线性化（写成 $\cos k\theta$ 的组合），并推出
$\displaystyle\int_0^{\pi/2}\cos^3\theta\,\dd\theta$.
\end{exercise}

\begin{exercise}[$\star\star$]\label{exo:g12:complex:7}
设点 $A$、$B$ 的对应复数分别为 $a = 1 + \iu$ 和 $b = 3 + 2\iu$，求
两点 $C$，使三角形 $ABC$ 为正三角形，它们是
$B$ 绕点 $A$ 分别旋转 $\pm\frac{\pi}{3}$ 所得的像。
\end{exercise}

\begin{exercise}[$\star\star\star$]\label{exo:g12:complex:8}
解方程 $z^4 = -4$（$\C$ 内），并画出其解。将
$z^4 + 4$ 分解为两个实系数二次多项式。
\end{exercise}

\begin{exercise}[$\star\star\star$]\label{exo:g12:complex:9}
对 $\theta \in \intoo{0}{2\pi}$ 与 $n \in \N$，计算
\[
S = 1 + \cos\theta + \cos 2\theta + \dots + \cos n\theta
\]
通过对等比数列
$\sum_{k=0}^n \eu^{\iu k\theta}$ 求和，并取实部。（答案：
$S = \dfrac{\sin\frac{(n+1)\theta}{2}}{\sin\frac{\theta}{2}}
\cos\frac{n\theta}{2}$.)
\end{exercise}

\begin{exercise}[$\star\star\star$]\label{exo:g12:complex:10}
设 $\omega = \eu^{2\iu\pi/5}$。
\begin{enumerate}
  \item 说明 $1 + \omega + \omega^2 + \omega^3 + \omega^4 = 0$ 成立。
  \item 设 $u = \omega + \omega^4$，$v = \omega^2 + \omega^3$。证明
        $u + v = -1$，$uv = -1$，并由此推出 $u = \frac{-1+\sqrt5}{2}$。
  \item 由此可得 $\cos\dfrac{2\pi}{5} = \dfrac{\sqrt5 - 1}{4}$。
\end{enumerate}
\end{exercise}

\section{问题：单位根镜中的正五边形}

\begin{problem}[{周末问题 --- 单位根的五次方根，
精确求出 $\cos 72^\circ$，黄金比为结果署名，
并揭示乘法实为旋转}]
\label{pb:g12:complex:1}
欧几里得能作出正五边形，但其中的缘由
\emph{为何}它能用尺规作出 --- 而区区的
$20^\circ$ 角则不能 --- 这一缘由两千年间始终
隐而未显。它藏于本章：五个五次单位根
（见\cref{thm:g12:complex:rootsofunity}）之和为零，而这一
行代数便从二次方程中解出 $\cos 72^\circ$，
方程 --- 仅含平方根，故可尺规作图，并且
自然以黄金比为其署名。围绕这一珍宝，有运算熟练、
棣莫弗的三角捷径，以及乘法被揭示为
几何。

\medskip
\noindent\textbf{第一部分 --- 运算熟练。}
\begin{enumerate}
  \item 计算 $(2 + \iu)(3 - \iu)$，\
        $\dfrac{1 + \iu}{1 - \iu}$
        （见\cref{met:g12:complex:division}），以及
        $\abs{3 + 4\iu}$.
  \item 化为指数形式：$1 + \iu$；\ $-2$；\
        $1 - \iu\sqrt3$.
  \item 利用指数形式计算 $(1 + \iu)^8$。
  \item 解方程 $z^2 = \iu$，再解 $z^2 - 2z + 5 = 0$
        （见\cref{thm:g12:complex:quadratic}）。
  \item 几何对照表
        （见\cref{prop:g12:complex:geometry}）：描述映射
        $z \mapsto \iu z$，并描述集合
        $\abs{z - (1 + \iu)} = 2$.
\end{enumerate}

\medskip
\noindent\textbf{第二部分 --- 五根，一个正五边形。}
设 $\omega = \eu^{2\iu\pi/5}$，考虑
$1, \omega, \omega^2, \omega^3, \omega^4$.
\begin{enumerate}[resume]
  \item 验证它们恰为方程 $z^5 = 1$ 的解，
        并描述它们在平面上所构成的图形。
  \item 证明它们的和为零。（提示：利用等比数列求和
        --- 或将和乘以 $1 - \omega$。）
  \item 取实部，并推出
        \[
        1 + 2\cos\frac{2\pi}{5} + 2\cos\frac{4\pi}{5} = 0 .
        \]
  \item 设 $c = \cos\frac{2\pi}{5}$，利用
        $\cos\frac{4\pi}{5}$ 的二倍角公式，将第 8 问
        化为关于 $c$ 的二次方程；解之，并
        得出
        \[
        \cos 72^\circ = \frac{\sqrt5 - 1}{4} .
        \]
  \item 数值验证之后，让黄金比为其结果署名：
        证明 $c = \frac{1}{2\varphi}$，其中
        $\varphi = \frac{1 + \sqrt5}{2}$
        （见\cref{pb:g10:algebra:1}）--- 并陈述著名的
        几何回响：在正五边形中，对角线是
        边长的 $\varphi$ 倍。
\end{enumerate}

\medskip
\noindent\textbf{第三部分 --- 棣莫弗的捷径。}
\begin{enumerate}[resume]
  \item 作图判定：$\cos 72^\circ$ 只需
        平方根（第 9 问），故正五边形
        可尺规作图；而 $\cos 20^\circ$
        满足一个不可约的三次方程
        （见\cref{pb:g12:trigo:1}），故 $60^\circ$ 角不能
        三等分。陈述初现的判别准则（平方根
        可作，立方根不可作），其完整理论 ---
        高斯 19 岁时的工作与正 17 边形 --- 载于
        大学阶段的分册之中。
  \item 用指数形式解方程 $z^4 = -4$，并复原
        \cref{exo:g12:complex:8}中的因式分解
        $z^4 + 4 = (z^2 - 2z + 2)(z^2 + 2z + 2)$，通过配对
        共轭复根。
  \item 用二项式定理展开 $(\cos\theta + \iu\sin\theta)^3$，
        结合棣莫弗公式，并
        在三行内重新推导 $\cos 3\theta = 4\cos^3\theta - 3\cos\theta$，
        与和角公式的途径相比较 ---
        见\cref{pb:g12:trigo:1}。
  \item 与\cref{exo:g12:complex:9}对照检验：当
        $\theta = \frac{2\pi}{5}$，$n = 4$ 时，和式
        $1 + \cos\theta + \dots + \cos 4\theta$ 等于多少，并
        它为何与第 8 问相容？
  \item 计算 $8\iu$ 的三个立方根，并用代数
        形式表示。
\end{enumerate}

\medskip
\noindent\textbf{第四部分 --- 乘法即几何。}
\begin{enumerate}[resume]
  \item 完整描述映射 $z \mapsto (1 + \iu)z$：
        它旋转的角度是多少，伸缩的倍数是多少，
        以及它对单位正方形有何作用？
  \item 利用共轭复数证明 $\abs{z_1 z_2} = \abs{z_1}\,\abs{z_2}$，
        （见\cref{prop:g12:complex:modulus}），再
        计算 $(1+\iu)^n$（$n = 1, 2, 3, 4$），并描述
        这些幂所描出的螺线。
  \item 设数 $\mathrm{j} = \eu^{2\iu\pi/3}$，证明
        $1 + \mathrm{j} + \mathrm{j}^2 = 0$，并在三角形
        $\left(1, \mathrm{j}, \mathrm{j}^2\right)$ 上验证经典判别法：
        $a + \mathrm{j} b + \mathrm{j}^2 c = 0$ 对
        某一定向的正三角形 $abc$ 成立。
  \item 代数基本定理（在本阶段不加证明地承认，
        由达朗贝尔--高斯给出）：每个多项式在
        $\C$ 上都能完全分解。由此推出实数域上的推论 ---
        每个实系数多项式都可分解为实系数的一次与二次因式的乘积，
        并指出第 12 问即是其中一例。
  \item 总结 --- $\C$ 的画像，每条一句：数
        化为点；乘法化为
        旋转加伸缩；单位根化为
        正多边形；棣莫弗把幂化为
        三角恒等式；代数封闭性保证：
        任何方程都不再需要更大的世界。尾声：
        第 10 问中相遇的是哪两个著名的数？
\end{enumerate}
\end{problem}
